% Written by paper/cpc/scripts/gen_tables.py.
\begin{table}[ht]
\centering
\small
\caption{Paired time ratios. The ratio is the readcon-core time divided by the peer time, so a value below 1 means readcon-core finished sooner. The interval is a bias-corrected and accelerated 95 percent interval of the median. Failures stay in the table and are not converted into times.}
\label{tab:ratios}
\footnotesize
\setlength{\tabcolsep}{4pt}
\begin{tabular}{llrrrr}
\toprule
Payload & Peer & OK & Fail & Ratio & Interval \\
\midrule
1 frame & ASE eon & 30 & 0 & 0.135 & [0.132, 0.137] \\
1 frame & chemfiles XYZ & 30 & 0 & 0.73 & [0.726, 0.733] \\
1 frame & h5py & 30 & 0 & 0.319 & [0.316, 0.321] \\
1 frame & MDAnalysis H5MD & 30 & 0 & 0.02 & [0.0197, 0.021] \\
100 frames & ASE eon & 30 & 0 & 0.178 & [0.178, 0.179] \\
100 frames & chemfiles XYZ & 30 & 0 & 1.11 & [1.11, 1.12] \\
100 frames & h5py & 30 & 0 & 11 & [11, 11] \\
100 frames & MDAnalysis H5MD & 30 & 0 & 0.0909 & [0.0897, 0.093] \\
10000 frames & ASE eon & 30 & 0 & 0.253 & [0.251, 0.256] \\
10000 frames & chemfiles XYZ & 30 & 0 & 1.58 & [1.57, 1.6] \\
10000 frames & h5py & 30 & 0 & 47.6 & [46.7, 48.5] \\
10000 frames & MDAnalysis H5MD & 30 & 0 & 0.16 & [0.157, 0.162] \\
Pt slab & ASE eon & 30 & 0 & 0.215 & [0.211, 0.218] \\
Pt slab & chemfiles XYZ & 30 & 0 & 1.04 & [1.03, 1.04] \\
Pt slab & h5py & 30 & 0 & 4.1 & [4.06, 4.13] \\
Pt slab & MDAnalysis H5MD & 30 & 0 & 0.365 & [0.359, 0.368] \\
\bottomrule
\end{tabular}
\end{table}
